\section{The manuscript, the workbench, and what the re-run verifies}\label{sec:setting}

\subsection{The imported theorem}

The manuscript's main theorem, which we call NS-00 as the workbench does, is the following \cite[Theorem~1.1]{OpenAINS}.

\begin{theorem}[NS-00; the manuscript's Theorem~1.1]\label{thm:ns00}
For every $\nu>0$ there are a force $f\in C_c^\infty(\R^3\times(0,\infty))$, a compact set $K$ and a smooth solution $(u,p)$ of the forced Navier--Stokes equations on $\R^3\times[0,1)$ with $u(\cdot,0)=0$, supports in $K$ and $\sup_t\|u(t)\|_2<\infty$, such that $\limsup_{t\uparrow1}\|u(t)\|_\infty=\infty$.
\end{theorem}
\status{Imported; \emph{not verified here}. Corollary~10.6 of the manuscript (src p.~125) is the periodic version. The workbench is a reconstruction and verification reader of this theorem, and its own ledger lists four items as unfinished.}

\subsection{The architecture of the proof}

The workbench's reader follows the manuscript's architecture, which we recall for orientation; none of these steps is verified here.
\begin{itemize}
\item \emph{Force as residual} (src p.~3). For any incompressible $(u,p)$ put $f:=\partial_tu+(u\cdot\nabla)u-\nu\Delta u+\nabla p$. The work is to make $f$ smooth through $t=1$ while $\|u\|_\infty\to\infty$.
\item \emph{A slender core} (src p.~4). It has radial scale $\ell_r\asymp\tau^{1/2}$, axial scale $\ell_z\asymp\tau^{1/2-h}$ and speed $|u_\theta|,|u_z|\asymp\tau^{-1/2-h}$, with $\tau=1-t$ and $0<h<1/100$. The core energy is $\asymp\tau^{1/2-3h}$ and the angular Reynolds number is $\mathrm{Re}_\theta\asymp\tau^{-h}$. The similarity coordinates are (src (3.2))
\[
\tau=q(1-\eta^2),\qquad z=q^D\eta,\qquad r^2=2qX,\qquad A=\tfrac12+h,\qquad D=\tfrac12-h.
\]
\item \emph{Profiles} (Theorem~4.6, src p.~32). The tangential residuals are exact radial divergences of a stress $q^{-A-1/2}T_0$, with $T_0=0$ on the inner region $X\le X_a$ and admissible on the annulus; the exterior is an exact heat solution.
\item \emph{Base correction} (Proposition~5.5, src p.~60). The residual of the base field equals stress divergences plus a flat remainder. The profile estimates (5.42) hold uniformly on $[X_{\mathrm{lo}},X_{\mathrm{hi}}]\times[-1,1]$ for the summed background, with the Stokes streamfunctions summed before differentiation.
\item \emph{Waves and corrections} (§§6--9). Pulses on an auxiliary torus with the matrix $J=\begin{psmallmatrix}3&1\\1&5\end{psmallmatrix}$ realise the stress through their covariance (Proposition~7.5); a four-step cycle raises the residual exponent $\sigma_j=\frac15+\frac j{10}\to\infty$; a summation with shrinking cutoffs (Proposition~9.9) gives the manuscript's Theorem~3.1.
\item \emph{The endpoint} (§10). Potential-first cutoffs (Proposition~10.1) are used; the force jets converge (Lemma~10.2) and are extended by a Borel series past $t=1$ (Lemma~10.3); energy and uniqueness follow (Lemmas~10.4--10.5). The growth path is (10.20)--(10.21), the viscosity map is (10.22), and the periodisation is Corollary~10.6.
\end{itemize}

\subsection{What the re-run of the workbench's scripts verifies}

The first-pass inventory classified the manuscript's statements and the thirteen reconstruction components of the workbench. Its ten scripts were re-run here with python-flint 0.9.0, sympy 1.14.0 and mpmath 1.3.0, and their outputs agree with the shipped ones line for line, apart from timing lines. They comprise $251$ identity checks and the $5$ conditions of the collision certificate, all passing, and $4$ negative controls, all correctly rejected. The workbench's own replay, of $9$ programs and $144$ result groups, passes as well.

The identities verified in this way are:
\begin{itemize}
\item the leading-field calculus, including $R^{(0)}=-(\partial_r+k/r)(q^{-A-1/2}T_0)$ in full;
\item the order-$n$ coefficient equations for $n\le2$, from the full residual;
\item the actual-shear correction (AS7)--(AS40) and the ambient projection;
\item the vacuum-hydrodynamics identities: the transverse-traceless lift, the Lichnerowicz reduction, the momentum constraint, the Kasner bijection, the Einstein check $\operatorname{Ric}=-4g/L^2$ and the Kretschmann limit;
\item all bridge identities of Section~\ref{sec:bridges};
\item the viscosity, parabolic and periodisation maps.
\end{itemize}
\status{Verified (re-run).}

Not verified here are NS-00 and its analytic estimates (the existence part of Theorem~4.6; Propositions 5.3, 5.5, 7.2--7.6, 9.6 and 9.9; Lemmas 10.2--10.5), the Volterra convergence and the contraction constants of the manuscript's Appendix~B, the stage-cycle exponent tables beyond the re-run of the bundle, and the continuation's boundary relation (2.3) with its Brown--York formulas (20)--(24).
